% Part: turing-machines
% Chapter: machines-computations
% Section: halting-states

\documentclass[../../../include/open-logic-section]{subfiles}

\begin{document}

% SEGMENT OLP-0259-S01

\olfileid{tur}{mac}{hal}
\olsection{நிற்றல் நிலைகள்}

\begin{explain}
நிறைவேற்றுவதற்குக் கட்டளை இல்லாதபோது மட்டுமே நமது பொறிகள் நிற்கும்
என்று நாம் வரையறுத்திருந்தாலும், டியூரிங் பொறிகளின் பொதுவான
குறிப்பீடுகளில் $h \in Q$ ஆக உள்ள, ஒதுக்கப்பட்ட ஒரு
\emph{நிற்றல் நிலை}~$h$ இருக்கும்.

ஒரு நிற்றல் நிலையின் அடிப்படைக் கருத்து எளியது: பொறி தனது செயல்பாட்டை
முடித்ததும் (அது உள்ளீட்டை ஏற்கத் தயாராக உள்ளது அல்லது வெளியீட்டை
எழுதி முடித்துள்ளது), தான் நிற்கும் நிலை~$h$ குச் செல்கிறது. சில
பொறிகளுக்கு இரு நிற்றல் நிலைகள் உள்ளன; ஒன்று உள்ளீட்டை ஏற்கவும்
மற்றொன்று உள்ளீட்டை மறுக்கவும் பயன்படும்.
\end{explain}

\begin{ex}\emph{நிற்றல் நிலைகள்}.
இக்கருத்தைத் தெளிவுபடுத்த, இரட்டை எண் பொறியில் ஒரு மாற்றத்துடன்
தொடங்குவோம். உள்ளீடு இரட்டை எண்ணுடையதாக இருந்தால் பொறியை நிலை~$q_0$
இல் நிற்கச் செய்வதற்குப் பதிலாக, பொறியை ஒரு நிற்றல் நிலைக்கு அனுப்பும்
ஒரு கட்டளையைச் சேர்க்கலாம்.
\[
\begin{tikzpicture}[->,>=stealth',shorten >=1pt,auto,node distance=2.8cm,
                    semithick]
  \tikzstyle{every state}=[fill=none,draw=black,text=black]

  \node[initial, state]         (A)                     {$q_0$};
  \node[state]         (B) [right of=A] {$q_1$};
  \node[state]         (C) [below of=A] {$h$};

  \path (A) edge [bend left] node {\TMtrans{\TMstroke}{\TMstroke}{R}} (B)
  	    edge node {\TMtrans{\TMblank}{\TMblank}{N}} (C)
        (B) edge [loop above] node {\TMtrans{\TMblank}{\TMblank}{R}} (B)
            edge [bend left] node {\TMtrans{\TMstroke}{\TMstroke}{R}} (A);
\end{tikzpicture}
\]

எடுத்துக்காட்டை மேலும் விரிவாக்குவோம். உள்ளீடு ஒற்றை எண்ணுடையது என்று
பொறி தீர்மானிக்கும்போது, அது ஒருபோதும் நிற்பதில்லை. கண்ணியில் இயங்கும்
கட்டளையை, மறுப்பு நிலை~$r$ குச் செல்லும் ஒரு கட்டளையால் மாற்றுவதன் மூலம்,
ஒரு \emph{மறுப்பு} நிலையைச் சேர்க்குமாறு பொறியை மாற்றலாம்.
\[
\begin{tikzpicture}[->,>=stealth',shorten >=1pt,auto,node distance=2.8cm,
                    semithick]
  \tikzstyle{every state}=[fill=none,draw=black,text=black]

  \node[initial,state]         (A)                     {$q_0$};
  \node[state]         (B) [right of=A] {$q_1$};
  \node[state]         (C) [below of=A] {$h$};
  \node[state]         (D) [below of=B] {$r$};

  \path (A) edge [bend left] node {\TMtrans{\TMstroke}{\TMstroke}{R}} (B)
  	    edge node {\TMtrans{\TMblank}{\TMblank}{N}} (C)
        (B) edge node {\TMtrans{\TMblank}{\TMblank}{N}} (D)
            edge [bend left] node {\TMtrans{\TMstroke}{\TMstroke}{R}} (A);
\end{tikzpicture}
\]
\end{ex}

\begin{explain}
இது போன்ற நேர்வுகளில் ஒதுக்கப்பட்ட ஒரு நிற்றல் நிலையைச் சேர்ப்பது
பயனளிக்கலாம்; ஏனெனில், குறிப்பிட்ட உள்ளீடுகளைப் பொறி எப்போது
ஏற்கிறது அல்லது மறுக்கிறது என்பதை அது வெளிப்படையாக்குகிறது. எனினும்,
ஒதுக்கப்பட்ட ஒரு நிற்றல் நிலையைச் சேர்ப்பதால் எந்தக் கணிப்புத் திறனும்
கூடுவதில்லை. இதேபோல், நிற்றல் பற்றிய குறைவான முறைசார் கருத்துக்கும்
அதற்கே உரிய பயன்கள் உள்ளன. இந்த அதிகாரத்தில் இதுவரை பயன்படுத்தப்பட்ட
நிற்றலின் வரையறை, \emph{நிற்றல் சிக்கலுக்கான} நிறுவலை உள்ளுணர்வுக்கு
இணங்கியதாகவும் எளிதில் எடுத்துக்காட்டக்கூடியதாகவும் ஆக்குகிறது.
இதனால், நமது மூல வரையறையையே தொடர்ந்து பயன்படுத்துகிறோம்.
\end{explain}

\end{document}
